\section{Software Documentation}\label{app:software}

This appendix provides installation instructions and API reference for the \textsc{ThinkingTrees} package. Full documentation is available at \url{https://github.com/thinking-trees/thinking-trees}.

\subsection{Repository Layout and Entry Points}

\paragraph{Package metadata.}
The Python package name is \texttt{thinking-trees} (see \texttt{pyproject.toml}).
Repository: \url{https://github.com/thinking-trees/thinking-trees}.
The import path used throughout the codebase is the repository-local \texttt{src/} package.

\paragraph{Primary entry points.}
\begin{itemize}[nosep]
    \item \texttt{main.py}: CLI for building a single OPS tree from a document.
    \item \texttt{src/training/run\_pipeline.py}: end-to-end training pipeline invoked by \texttt{scripts/run\_training\_pipeline.sh}.
    \item \texttt{scripts/start\_dual\_servers.sh} and \texttt{scripts/start\_vllm.sh}: vLLM server launchers driven by \texttt{config/settings.yaml}.
\end{itemize}

\paragraph{Core modules.}
\begin{itemize}[nosep]
    \item \texttt{src/core/}: LLM clients, batching, scoring/oracles, protocols, and request orchestration.
    \item \texttt{src/tree/}: tree construction, auditing, verification, and labeled datasets.
    \item \texttt{src/training/}: DSPy optimization, preference collection, judges, metrics, and training loops.
    \item \texttt{src/tasks/} and \texttt{src/datasets/}: pluggable task and dataset registries.
    \item \texttt{src/pipelines/}: batched multi-document pipeline for high-throughput inference.
\end{itemize}

\subsection{Installation}

\paragraph{Requirements.}
\begin{itemize}[nosep]
    \item Python 3.10+
    \item \texttt{dspy>=3.0.0}
    \item \texttt{openai>=1.0.0} (for OpenAI-compatible endpoints)
    \item Optional: \texttt{vllm>=0.12.0} for local inference
\end{itemize}

\paragraph{Basic Installation.}
\begin{lstlisting}
pip install thinking-trees
\end{lstlisting}

\paragraph{Development Installation.}
\begin{lstlisting}
git clone https://github.com/thinking-trees/thinking-trees.git
cd ThinkingTrees
pip install -e ".[dev]"
\end{lstlisting}

\paragraph{Optional Dependencies.}
\begin{lstlisting}
pip install thinking-trees[vllm]      # vLLM backend
pip install thinking-trees[manifesto] # Manifesto dataset utilities
pip install thinking-trees[trl]       # TRL training
pip install thinking-trees[all]       # All optional deps
\end{lstlisting}

\subsection{Quick Start}

\paragraph{Building a Tree.}
\begin{lstlisting}
import asyncio
from src.core.batch_processor import AsyncBatchLLMClient
from src.core.strategy import BatchedStrategy
from src.tree.builder import TreeBuilder, BuildConfig

async def build_tree(text: str, rubric: str) -> None:
    async with AsyncBatchLLMClient(base_url="http://localhost:8000/v1") as client:
        strategy = BatchedStrategy(client, max_tokens=500, unified_mode=True)
        builder = TreeBuilder(strategy, config=BuildConfig(max_chunk_chars=2000))
        result = await builder.build(text, rubric)
        print(f"Root summary: {result.tree.final_summary[:200]}...")

asyncio.run(build_tree(document, rubric="Preserve key information and entities."))
\end{lstlisting}

\paragraph{Auditing a Tree.}
\begin{lstlisting}
from src.tree.auditor import Auditor, AuditConfig, SimpleScorer

scorer = SimpleScorer()
auditor = Auditor(oracle=scorer, config=AuditConfig(sample_budget=25))
report = auditor.audit_tree(tree)

print(f"Nodes audited: {report.nodes_audited}")
print(f"Sufficiency rate: {report.sufficiency_rate:.4f}")
print(f"Merge rate: {report.merge_rate:.4f}")
\end{lstlisting}

\paragraph{Collecting Preferences.}
\begin{lstlisting}
from src.training.preference.types import PreferenceDataset
from src.training.preference.collector import PreferenceCollector

collector = PreferenceCollector(summarizer=my_dspy_summarizer, k=4)
pairs = collector.collect_pairs_for_example(
    example_id="doc1",
    original_text=document,
    rubric="Preserve key information",
    law_type="sufficiency",
)

# Export for DPO training
dataset = PreferenceDataset(pairs)
dataset.save("preferences.json")
\end{lstlisting}

\subsection{API Reference}

\subsubsection{Core Data Models}

\paragraph{\texttt{Node}.}
\begin{lstlisting}
@dataclass
class Node:
    id: str
    level: int
    raw_text_span: Optional[str]
    summary: str
    left_child: Optional[Node]
    right_child: Optional[Node]
    parent: Optional[Node]
    audit_result: AuditResult

    @property
    def is_leaf(self) -> bool: ...
\end{lstlisting}

\paragraph{\texttt{Tree}.}
\begin{lstlisting}
@dataclass
class Tree:
    root: Node
    rubric: str = ""
    metadata: dict = {}

    def traverse_preorder(self) -> Iterator[Node]: ...
    def save(self, path: Path) -> None: ...
    @classmethod
    def load(cls, path: Path) -> "Tree": ...
\end{lstlisting}

\paragraph{\texttt{DocumentSample} and \texttt{DocumentResult}.}
\begin{lstlisting}
@dataclass
class DocumentSample:
    doc_id: str
    text: str
    reference_score: Optional[float] = None
    metadata: Dict[str, Any] = {}

@dataclass
class DocumentResult:
    doc_id: str
    final_summary: str
    reference_score: Optional[float]
    estimated_score: Optional[float]
    tree_height: Optional[int]
    tree_nodes: Optional[int]
    error: Optional[str]
\end{lstlisting}

\subsubsection{TreeBuilder}

\begin{lstlisting}
@dataclass
class BuildConfig:
    max_chunk_chars: int = 2000
    min_chunk_chars: int = 100
    chunk_strategy: str = "sentence"
    merge_strategy: str = "binary"
    pipelined: bool = True

class TreeBuilder:
    def __init__(self, strategy: SummarizationStrategy, config: BuildConfig = None): ...
    async def build(self, text: str, rubric: str = "") -> BuildResult: ...
    async def build_from_file(self, filepath: Path, rubric: str = "") -> BuildResult: ...
    def build_sync(self, text: str, rubric: str = "") -> BuildResult: ...

@dataclass
class BuildResult:
    tree: Tree
    chunks_created: int
    nodes_created: int
    levels_created: int
    errors: List[str]
    preferences: List[PreferencePair]
\end{lstlisting}

\subsubsection{Auditing and Checks}

\begin{lstlisting}
@dataclass
class AuditConfig:
    sample_budget: int = 10
    discrepancy_threshold: float = 0.1
    audit_leaves: bool = True
    audit_internal: bool = True
    audit_idempotence: bool = True
    audit_substitution: bool = True

class Auditor:
    def __init__(self, oracle: ScoringOracle, config: AuditConfig = None): ...
    def audit_tree(self, tree: Tree) -> AuditReport: ...
\end{lstlisting}

\subsubsection{Summarization Strategies}

\paragraph{\texttt{BatchedStrategy}.}
High-throughput batched inference:
\begin{lstlisting}
strategy = BatchedStrategy(client, max_tokens=500, unified_mode=True)
\end{lstlisting}

\paragraph{\texttt{DSPyStrategy}.}
Optimizable with DSPy:
\begin{lstlisting}
strategy = DSPyStrategy(
    leaf_module=GenericSummarizer(),
    merge_module=GenericMerger(),
)
\end{lstlisting}

\subsubsection{Task Plugin Interface}

\begin{lstlisting}
class AbstractTask(ABC):
    @property
    def name(self) -> str: ...
    @property
    def scale(self) -> Optional[ScaleDefinition]: ...
    def create_rubric(self) -> str: ...
    def create_predictor(self, retriever=None, config=None) -> dspy.Module: ...
    def create_prompt_builders(self) -> PromptBuilders: ...
    def create_metric(self, with_feedback: bool = True) -> Callable: ...

task = ScoringTask(
    name="my_task",
    scale=ScaleDefinition(name="score", min_value=-100, max_value=100),
    rubric="Preserve key information for the task...",
)
\end{lstlisting}

\subsection{LLM Clients and Batching}

\paragraph{Single-request client.}
\texttt{src/core/llm\_client.py} provides \texttt{LLMConfig} and \texttt{LLMClient}
for OpenAI-compatible endpoints (vLLM, SGLang, OpenAI API). It supports:
retry limits, request timeouts, and optional LRU response caching.

\paragraph{Batched client.}
\texttt{AsyncBatchLLMClient} pools requests, sends batched HTTP calls to vLLM,
tracks throughput statistics, and returns structured \texttt{BatchResponse} objects.
Global scheduling across documents is implemented in \texttt{BatchTreeOrchestrator}.
Concurrency defaults are centralized in \texttt{src/config/concurrency.py}.

\subsection{Configuration Reference}

Configuration is YAML-based with defaults in \texttt{config/settings.yaml}.
Additional runtime configs include \texttt{config/inference.yaml},
\texttt{config/training.yaml}, and \texttt{config/audit.yaml}.

\paragraph{Key settings.}
\begin{itemize}[nosep]
    \item \texttt{servers.task\_model\_url}, \texttt{servers.genrm\_url} (override with \texttt{TASK\_MODEL\_URL}, \texttt{GENRM\_URL})
    \item \texttt{tasks.default}, \texttt{datasets.default} (override with \texttt{TASK}, \texttt{DATASET})
    \item \texttt{generation.*} for summarizer/oracle/judge sampling parameters
    \item \texttt{chunking.max\_tokens} (settings) or \texttt{chunking.max\_chars} (inference) plus \texttt{context\_window.*} for input/output allocation
    \item \texttt{audit.*} for sample budgets and thresholds
    \item \texttt{optimization.*} for DSPy optimization control
\end{itemize}

\paragraph{Environment overrides for LLM client.}
LLM client defaults can be overridden via \texttt{LLM\_BASE\_URL},
\texttt{LLM\_MODEL}, and \texttt{LLM\_API\_KEY}.

\subsection{Dataset and Artifact Formats}

\paragraph{JSONL dataset.}
Each line is a JSON object with required \texttt{text} (or \texttt{content})
and optional \texttt{doc\_id}/\texttt{id}. Optional numeric score fields:
\texttt{reference\_score} or \texttt{score} (\texttt{src/datasets/jsonl.py}).

\paragraph{Manifesto dataset.}
\texttt{src/datasets/manifesto.py} wraps the Manifesto Project loader and
maps fields like \texttt{party\_name}, \texttt{country\_code}, and \texttt{year}
into \texttt{DocumentSample.metadata}.

\paragraph{Tree serialization.}
\texttt{Tree.save()} writes a JSON blob with \texttt{root\_id}, \texttt{rubric},
\texttt{metadata}, and a flat \texttt{nodes} list (see \texttt{Tree.to\_dict}).

\paragraph{Preference datasets.}
Preference pairs are stored as JSON with fields from \texttt{PreferencePair}
(see \texttt{src/training/preference/types.py}). The training pipeline writes:
\texttt{preferences/preferences\_ops\_tree\_*.json},
\texttt{preferences/dpo\_ops\_tree\_*.json}, and
\texttt{preferences/tree\_stats\_*.json}.

\subsection{Scripted Workflows}

\paragraph{Single-document CLI.}
\begin{lstlisting}
python main.py --input data/raw/example.txt --port 8000 --output outputs/tree.json
\end{lstlisting}

\paragraph{Server control.}
\begin{lstlisting}
./scripts/start_dual_servers.sh
./scripts/start_dual_servers.sh --small-only
./scripts/start_dual_servers.sh --large-only
./scripts/start_vllm.sh nemotron-30b-fp8
./scripts/stop_small_servers.sh --all
\end{lstlisting}

\paragraph{Training pipeline.}
\begin{lstlisting}
./scripts/run_training_pipeline.sh \
  --output-dir outputs/train_$(date +%Y%m%d_%H%M) \
  --train-samples 100 --val-samples 30 --test-samples 30 \
  --optimizer bootstrap_random_search --optimizer-budget heavy \
  --n-iterations 2 --enable-genrm --genrm-port 8001
\end{lstlisting}

\subsection{Troubleshooting}

\paragraph{Rate limiting.}
Use \texttt{AsyncBatchLLMClient} to control concurrency and batching:
\begin{lstlisting}
client = AsyncBatchLLMClient(
    base_url="http://localhost:8000/v1",
    max_concurrent=50,
    batch_size=20,
    batch_timeout=0.05,
)
\end{lstlisting}

\paragraph{Memory issues.}
For large documents, reduce request sizes or chunk sizes:
\begin{lstlisting}
config = BuildConfig(max_chunk_chars=1500)
\end{lstlisting}

\paragraph{Timeout errors.}
Increase request timeouts for slow models:
\begin{lstlisting}
client = AsyncBatchLLMClient(request_timeout=600.0)
\end{lstlisting}
